% 本页逐项据用户补拍的原书第68页重录，不沿用旧稿误配题目。
\QHeading{第1题　口算。}
\begin{oralgrid}
\OralItem{1}{0.2+0.3}{0.5}
\OralItem{2}{0.5+0.4}{0.9}
\OralItem{3}{1.25+0.55}{1.8}
\OralItem{4}{2.3+1.1}{3.4}
\OralItem{5}{0.8+0.7}{1.5}
\OralItem{6}{3.65+0.6}{4.25}
\OralItem{7}{5.5+2.5}{8}
\OralItem{8}{0.36+0.4}{0.76}
\OralItem{9}{1.25+0.75}{2}
\OralItem{10}{0.9-0.2}{0.7}
\OralItem{11}{1.05-0.8}{0.25}
\OralItem{12}{2.6-1.3}{1.3}
\OralItem{13}{3.25-1.5}{1.75}
\OralItem{14}{4.8-2.8}{2}
\OralItem{15}{1.25-0.6}{0.65}
\OralItem{16}{2.3+3.1}{5.4}
\end{oralgrid}
\QHeading{第2题　脱式计算，能简算的要简算。}
\CalcBlock[方法一：简便计算]{1}{\begin{aligned}[t]&\phantom{={}}3.7+2.88+6.3\\&\color{answerred}=3.7+6.3+2.88\\&\color{answerred}=(3.7+6.3)+2.88\\&\color{answerred}=10+2.88\\&\color{answerred}=12.88\end{aligned}}
\MethodBlock{方法二：按运算顺序核对}{与上面的简便算法比较，得数相同。}{\begin{aligned}[t]&\phantom{={}}3.7+2.88+6.3\\&=6.58+6.3\\&=12.88\end{aligned}}
\CalcBlock[方法一：简便计算]{2}{\begin{aligned}[t]&\phantom{={}}5.35+1.9+4.65\\&\color{answerred}=5.35+4.65+1.9\\&\color{answerred}=(5.35+4.65)+1.9\\&\color{answerred}=10+1.9\\&\color{answerred}=11.9\end{aligned}}
\MethodBlock{方法二：按运算顺序核对}{与上面的简便算法比较，得数相同。}{\begin{aligned}[t]&\phantom{={}}5.35+1.9+4.65\\&=7.25+4.65\\&=11.9\end{aligned}}
\CalcBlock[方法一：简便计算]{3}{\begin{aligned}[t]&\phantom{={}}12.8+4.55+5.45\\&\color{answerred}=12.8+(4.55+5.45)\\&\color{answerred}=12.8+10\\&\color{answerred}=22.8\end{aligned}}
\MethodBlock{方法二：按运算顺序核对}{与上面的简便算法比较，得数相同。}{\begin{aligned}[t]&\phantom{={}}12.8+4.55+5.45\\&=17.35+5.45\\&=22.8\end{aligned}}
\CalcBlock[方法一：简便计算]{4}{\begin{aligned}[t]&\phantom{={}}20.6-3.2-6.8\\&\color{answerred}=20.6-(3.2+6.8)\\&\color{answerred}=20.6-10\\&\color{answerred}=10.6\end{aligned}}
\MethodBlock{方法二：按运算顺序核对}{与上面的简便算法比较，得数相同。}{\begin{aligned}[t]&\phantom{={}}20.6-3.2-6.8\\&=17.4-6.8\\&=10.6\end{aligned}}
\CalcBlock[方法一：简便计算]{5}{\begin{aligned}[t]&\phantom{={}}15.9-2.48-7.52\\&\color{answerred}=15.9-(2.48+7.52)\\&\color{answerred}=15.9-10\\&\color{answerred}=5.9\end{aligned}}
\MethodBlock{方法二：按运算顺序核对}{与上面的简便算法比较，得数相同。}{\begin{aligned}[t]&\phantom{={}}15.9-2.48-7.52\\&=13.42-7.52\\&=5.9\end{aligned}}
\CalcBlock[方法一：简便计算]{6}{\begin{aligned}[t]&\phantom{={}}8.36+(1.64+5.7)\\&\color{answerred}=(8.36+1.64)+5.7\\&\color{answerred}=10+5.7\\&\color{answerred}=15.7\end{aligned}}
\MethodBlock{方法二：按运算顺序核对}{与上面的简便算法比较，得数相同。}{\begin{aligned}[t]&\phantom{={}}8.36+(1.64+5.7)\\&=8.36+7.34\\&=15.7\end{aligned}}
\QHeading{第3题　解决问题。}
\Problem{① 一支钢笔售价8.5元，一本笔记本售价3.8元，一块橡皮售价1.2元。妈妈买了一支钢笔、一本笔记本和一块橡皮，一共需要多少钱？}
\MethodTitle{方法一：先把笔记本和橡皮的价钱凑整}
\Answer{\Formula{\begin{aligned}[t]&\phantom{={}}8.5+3.8+1.2\\&=8.5+(3.8+1.2)\\&=8.5+5\\&=13.5\text{（元）}\end{aligned}}
答：一共需要13.5元。}
\MethodTitle{方法二：先算钢笔和笔记本的价钱}
\Answer{\Formula{\begin{aligned}[t]&\phantom{={}}8.5+3.8+1.2\\&=12.3+1.2\\&=13.5\text{（元）}\end{aligned}}
答：一共需要13.5元。}
\Problem{② 妈妈的电子钱包里有50.8元，买菜花了18.65元，买水果又花了21.35元。妈妈的电子钱包里还剩下多少钱？}
\MethodTitle{方法一：先求一共花了多少钱}
\Answer{\Formula{\begin{aligned}[t]&\phantom{={}}50.8-18.65-21.35\\&=50.8-(18.65+21.35)\\&=50.8-40\\&=10.8\text{（元）}\end{aligned}}
答：妈妈的电子钱包里还剩下10.8元。}
\MethodTitle{方法二：依次减去两笔支出}
\Answer{\Formula{\begin{aligned}[t]&\phantom{={}}50.8-18.65-21.35\\&=32.15-21.35\\&=10.8\text{（元）}\end{aligned}}
答：妈妈的电子钱包里还剩下10.8元。}
\MethodTitle{方法三：把买菜钱按20元计算，再补偿}
\Answer{\Formula{\begin{aligned}[t]&\phantom{={}}50.8-18.65-21.35\\&=50.8-20+1.35-21.35\\&=30.8+1.35-21.35\\&=30.8-(21.35-1.35)\\&=30.8-20\\&=10.8\text{（元）}\end{aligned}}
答：妈妈的电子钱包里还剩下10.8元。}
